% 来源：OEBPS/Text/01_27.xhtml（OPF spine 顺序；仅第一册）
\cleardoublepage
\BookEntry[伯纳德·曼德维尔《蜜蜂的寓言》]{伯纳德·曼德维尔《蜜蜂的寓言》}

\BookCover{images/书单13蜜蜂的寓言.jpg}{[荷] 伯纳德·曼德维尔著| 肖聿译| 中国社会科学出版社2002年版}

\begin{BookSummaryBox}
\textbf{“用英国语言写出的最邪恶、最聪明的书”}
\par
\end{BookSummaryBox}

\subsubsection*{这本书谈什么？}

《蜜蜂的寓言》一书的核心，是道出了西方思想史上著名的曼德维尔悖论，即“私人恶德即公共利益”。从“人是自私自利的”这一观点出发，曼德维尔认为，在人品美德和商业社会之间存在着根本张力。因此，在他看来，所谓的试图以道德说教为手段，以“公共精神”为基础，建立一个美好繁荣的社会是不现实的。因为“不是人之美德和良善，而是制度性规则使道德品行不好的人也能为社会的福祉服务；而制约人们自私的恶行的绝大多数制度和规则，都是在充满罪恶和欺骗的商业行为的人类实践中产生”。
\par

曼德维尔的思想引起许多争议，该书18世纪在英国出版后，曾引起轩然大波，招致欧洲思想界乃至社会各界的批判。但尽管如此，这本书在西方思想史上的影响却还是不可小觑，影响了休谟、亚当·斯密和哈耶克等大师。
\par

\subsubsection*{芒格为什么谈起这本书？}

在芒格“论学院派经济学”的演讲中，他谈到经济学中总有一些无法避免的悖论。比如，恶行效应引发的一个问题与妒忌有关。他说：
\par

\begin{quotation}
妒忌在摩西的律法中是饱受指责的。你们可能记得希伯来人刻在墓碑上的文字：你们不能觊觎邻居的驴子，你们不能垂涎邻居的女仆，你们不能贪图……这些古代犹太人知道妒忌的人们是什么德性，也知道这些人会惹出什么麻烦。他们对妒忌真的绝不姑息，他们这么做是正确的。但曼德维尔令人信服地——反正我是信服的——证明妒忌会极大地推动消费意欲。所以呢，妒忌既是摩西十诫严厉禁止的糟糕恶行，却又是促使经济增长的驱动力。
\par
\end{quotation}

芒格提到的曼德维尔，是哲学家和讽刺作家，他在1705年发表了一首讽刺时局的诗歌：《蜜蜂的寓言，或，私人的恶行，公共的利益》。曼德维尔的哲学观是，利他主义损害了国家利益和学术进步，人类的自私心反倒是真正推动进步的引擎。因而他得到了“私心即是公益”的悖论。
\par

\subsubsection*{《蜜蜂的寓言》作者伯纳德·曼德维尔如是说}

“私人恶德即公众利益。”
\par

“若要界定作为奖赏的、表现得最充分的光荣，要廓清对光荣最多见的说法，便应当说：光荣乃是一种至高无上的快乐。一个人意识到自己做出了高尚的行为，便会在自爱当中享受这种乐，而他亦同时期待来自他人的喝彩。”
\par

“有理智、有教养者以最为娴熟的技巧掩盖自傲之时，亦是他最能取悦于自傲之心的时刻。”
\par
